\begin{itemize}
    \item Introduzione: capitolo non numerato, va all'inizio. Descrive la situazione "iniziale", quella di quando hai cominciato lo stage, descrive obiettivi e metodologie, e fornisce una panoramica della relazione finale. 
    \item La Ciccia: 2-3 capitoli dove spieghi cosa hai fatto nello stage. Nel caso di relazioni teoriche sono spesso 
    \begin{itemize}
        \item Introduzione teorica
        \item Studio (elaborazione)
        \item Implementazione se c'è
    \end{itemize}
    In ogni caso di solito il primo capitolo è sempre di nozioni utili a comprendere il resto della tesi. Un po' come se fosse un formulario e un libro di testo insieme, ma condensato in 5 pagine.
    \item Conclusioni: capitolo numerato che ripercorre la relazione mostrando i risultati ottenuti e gli sviluppi futuri che ne potrebbero derivare.
\end{itemize}

%\textbf{OGNI. CAPITOLO. DEVE. AVERE. UN. SUO. FILE. LATEX.}




% Contatore per le note
\newcounter{NoteCounter}
% Titolo per le note
\newcommand{\noteTitle}{\color{Green}\vspace{10pt}\noindent\textbf{\refstepcounter{NoteCounter}Nota n.\theNoteCounter }}

% Comando per inserire una nota
\newcommand{\nota}[1]{%
\ifnum0=\noteEnabled\relax
\else
    \noteTitle{}
    \textit{#1}
    \color{black}
\fi
}



\begin{figure}[h]
  \centering
  \includegraphics[trim={2cm 0cm 0cm 0cm}, scale=0.25, keepaspectratio]
  {Immagini/Dal 9 maggio.png}
  \caption{Area chart.}
\end{figure}


\begin{figure}[h]
    \centering
    \begin{subfigure}[b]{0.4\textwidth}
        \centering
        \includegraphics[width=7cm]{Immagini/Capitolo2/privateKeyCriptography.PNG}
        \caption{Crittografia simmetrica}
        \label{simmetric}
    \end{subfigure}
    \hspace{0.5cm}
    \begin{subfigure}[b]{0.4\textwidth}
        \centering
        \includegraphics[width=7cm]{Immagini/Capitolo2/publicKeyCriptography.PNG}
        \caption{Crittografia asimmetrica}
        \label{asimmetric}
    \end{subfigure}
    \caption{Confronto tra i due principali schemi crittografici~\cite{hankerson}.}
    \label{fig:confronto}
\end{figure}
\vspace{2mm}